% Mirrors `AlgoLib.Algorithms.Graph.Traversal.*`
%
\subsection{Breadth-first search}
% Mirrors `AlgoLib.Algorithms.Graph.Traversal.BFS`

Throughout, $G$ is a simple directed graph whose vertex set $\V{G}$ is finite and given as
data, and whose arc membership is decidable. Undirected graphs are searched through their
symmetric orientation (Definition~\ref{def:symmetric-orientation}), so one search serves
both kinds of graph.

\begin{definition}[BFS layers]
  \label{def:bfs-layers}
  \lean{AlgoLib.SimpleDiGraph.bfsLayer, AlgoLib.SimpleDiGraph.bfsVisited}
  \leanok
  Let $s \in \V{G}$. The \emph{layers} of breadth-first search from $s$ are the sets
  \[
    L_0 = \{s\}, \qquad
    L_{k+1} = N^+(L_k) \setminus (L_0 \cup \dots \cup L_k),
  \]
  where $N^+(L)$ denotes the set of out-neighbours of the vertices of $L$. The set of
  vertices \emph{visited} within $k$ rounds is $V_k = L_0 \cup \dots \cup L_k$. When
  $s \notin \V{G}$ every layer is empty.
  In Lean the search is a state $(V_k, L_k)$ iterated $k$ times from $(\{s\}, \{s\})$,
  so that the definition reduces and \texttt{decide} can evaluate it.
\end{definition}

\begin{lemma}[Soundness]
  \label{lem:bfs-sound}
  \lean{AlgoLib.SimpleDiGraph.exists_walk_of_mem_bfsLayer}
  \leanok
  \uses{def:bfs-layers}
  Every vertex of $L_k$ is the endpoint of a walk of length $k$ from $s$.
\end{lemma}
\begin{proof}
  \leanok
  By induction on $k$: a vertex of $L_{k+1}$ is an out-neighbour of a vertex of $L_k$,
  and appending that arc to the walk given by the induction hypothesis yields a walk of
  length $k + 1$.
\end{proof}

\begin{lemma}[Completeness]
  \label{lem:bfs-complete}
  \lean{AlgoLib.SimpleDiGraph.mem_bfsVisited_of_isVertexSeqIn}
  \leanok
  \uses{def:bfs-layers}
  The endpoint of a walk of length $\ell$ from $s$ lies in $V_\ell$.
\end{lemma}
\begin{proof}
  \leanok
  By induction on the walk. The key step is that the out-neighbours of a vertex of $V_k$
  lie in $V_{k+1}$: such a vertex lies in some $L_j$ with $j \le k$, and its
  out-neighbours are either already in $V_j$ or in $L_{j+1}$.
\end{proof}

\begin{lemma}[Termination]
  \label{lem:bfs-termination}
  \lean{AlgoLib.SimpleDiGraph.lt_card_of_nonempty_bfsLayer}
  \leanok
  \uses{def:bfs-layers}
  If $L_k \neq \emptyset$ then $k < |\V{G}|$. Hence the search from $s$ has exhausted every
  layer after $|\V{G}|$ rounds.
\end{lemma}
\begin{proof}
  \leanok
  An empty layer is followed by an empty layer, so $L_0, \dots, L_k$ are all non-empty.
  They are pairwise disjoint subsets of $\V{G}$, so $k + 1 \le |V_k| \le |\V{G}|$.
\end{proof}

\begin{theorem}[The layers are the spheres]
  \label{thm:bfs-layers-spheres}
  \lean{AlgoLib.SimpleDiGraph.mem_bfsLayer_iff_dist_eq}
  \leanok
  \uses{def:bfs-layers, def:dist}
  For every vertex $v$ and every $k \in \N$,
  \[
    v \in L_k \iff \dist(s, v) = k .
  \]
\end{theorem}
\begin{proof}
  \leanok
  \uses{lem:bfs-sound, lem:bfs-complete}
  Let $v \in L_k$. Soundness gives $\dist(s, v) \le k$. For the reverse inequality there
  is nothing to show when $k = 0$; when $k = j + 1$, a walk from $s$ to $v$ of length at
  most $j$ would place $v$ in $V_j$ by completeness, whereas $L_{j+1}$ avoids $V_j$ by
  construction, so $\dist(s, v) > j$, i.e. $\dist(s, v) \ge k$. Conversely, if
  $\dist(s, v) = k$ then a shortest walk places $v \in V_k$, so $v \in L_j$ for some
  $j \le k$, and the first part gives $j = k$.
\end{proof}

\begin{theorem}[BFS computes reachability and distance]
  \label{thm:bfs-correct}
  \lean{AlgoLib.SimpleDiGraph.mem_bfsReachableFinset_iff, AlgoLib.SimpleDiGraph.bfsDist_eq_dist}
  \leanok
  \uses{def:bfs-layers, def:reachable-directed, def:dist}
  The set $V_{|\V{G}|}$ is exactly the set of vertices reachable from $s$, and the least
  $k < |\V{G}|$ with $v \in L_k$ (or $\infty$ if there is none) is exactly $\dist(s, v)$.
\end{theorem}
\begin{proof}
  \leanok
  \uses{lem:bfs-sound, lem:bfs-complete, lem:bfs-termination}
  Soundness gives one inclusion and one inequality; completeness together with the
  termination bound gives the other.
\end{proof}
